\documentclass[11pt]{article}
\usepackage[a4paper,margin=23mm]{geometry}
\usepackage[T1]{fontenc}
\usepackage{lmodern,amsmath,amssymb,amsthm,mathtools,microtype,booktabs}
\usepackage[colorlinks=true,linkcolor=blue,citecolor=blue,urlcolor=blue]{hyperref}
\newcommand{\EE}{\mathbb E}
\newcommand{\tr}{\operatorname{tr}}
\newcommand{\LL}{\mathcal L}
\newcommand{\NN}{\mathcal N}
\newcommand{\WW}{\mathcal W}
\newcommand{\PP}{\mathcal P_H}
\newcommand{\BB}{\mathcal B}
\newcommand{\FF}{\mathcal F}
\newcommand{\ad}{\operatorname{ad}}
\newtheorem{theorem}{Theorem}
\title{Summable fourth-norm tails for the complete singlet bilocal\\
\large Common-disorder finite-time four-point functions in complex DSSYK}
\author{Alan Day \and Chengheng\thanks{AI research assistant (OpenAI).}}
\date{7 November 2026\quad Local research draft}
\begin{document}\maketitle
\begin{abstract}
We close the finite-time fourth-norm tail left open by the preceding
bilocal reconstruction. Restricted Gaussian covariances and an ancilla
contraction bound control pure marked words at every degree. An exact
sequence of completely positive charge-transfer maps then controls the
actual neutral covariance residual. Clifford Riesz inequalities,
charge finite differences and an inverse Jacobi expansion give the
corresponding bound for ordinary charged powers. The resulting tails
are uniform at finite size. Together with the previously established
same-disorder strong replacements and return estimate, they construct
the full finite-time bilocal four-point function by convergent finite
pairing sums. The result concerns the specified microscopic model and
tracial state; the state/polarization and bulk-generator dictionary are
the next physical questions.
\end{abstract}

\section{Objects and result}
Use standard CAR, $\{c_i,c_j^\dagger\}=\delta_{ij}$, and
\begin{equation}
 H_N=\sqrt{w_N}\sum_{|I|=|J|=p}J_{IJ}C_I^\dagger C_J,
 \qquad w_N=\frac{4^pJ_s^2}{\lambda_N}\left(\frac{p!}{N^p}\right)^2.
\end{equation}
The Gaussian matrix $J$ is Hermitian, with
$\EE J_{IJ}J_{KL}=\delta_{IL}\delta_{JK}$; $I,J$ may overlap.
Every sector uses this same $J$. Set
$\tau_N=\EE_J2^{-N}\operatorname{Tr}$ and
$\|X\|_r=(\tau_N|X|^r)^{1/r}$. The state is $\beta=\mu=0$.
Take even $N$, $p\to\infty$, $\lambda_N=p^2/N\to\lambda>0$
in a positive compact set, fixed $J_s$, $\hbar=1$, and string time.
All norm orders and real time windows below are fixed.

For neutral operators define $D=Q-\sum_i c_i^\dagger(\cdot)c_i$
and the half-Clifford length $\LL\Gamma_A=|A|\Gamma_A/2$.
$\NN$ and $\WW_j$ instead denote coupling Gaussian number and chaos
projection. Write $Z_j=\WW_j(H^j)$ and $C=J_s^2/\lambda$.
The actual bilocal is the previously defined
$\BB_N(t)=(1-\PP)m e^{-iH_mt/2}F_m(t)e^{-iH_mt/2}$,
where $F_m(t)=m^{-1}\sum_i c_i^\dagger e^{itH_{m-1}}c_i$;
inactive sectors have zero dynamic field. Pinching $\PP$ retains
whole degenerate energy blocks.

\newpage
\begin{theorem}[Finite-time four-point construction]\label{thm}
The charged and neutral series in the two-component memory reduction
have uniformly summable tails in every fixed finite $S_r$ norm,
$r\ge2$. In particular their
degree-$n$ building blocks admit bounds of the form
\begin{equation}\label{envelope}
 \|X_{n,N}\|_r\le K_r(1+n)^{b_r}L_r^n\sqrt{n!},
\end{equation}
with constants independent of $N,n$. Consequently the joint ordered
four-point functions of $\BB_N(t)$ and its adjoint, with arbitrary
fixed center-time conjugations, converge locally uniformly in all
these real times. Their limit is the convergent pairing construction
in Section~\ref{construction}.
\end{theorem}
The complete compressed-bath memory is connected by an $S_r$ strong
approximation; no separate claim about its own Taylor tail is needed.
This is a model theorem assembled from the new estimates below and
the explicit preceding strong interfaces \cite{Four,Joint,Recon}.
It does not assert a uniform growing-time limit or interchange the
double-scaling limit with $\lambda\downarrow0$.

\section{Restricted covariances without absolute CAR weights}
Let $A_{IJ}=\sqrt{w_N}C_I^\dagger C_J$. On charge $m$,
\begin{equation}
 \EE H_m^2=C_mI,\qquad
 C_m=w_N\binom mp\binom{N-m+p}p\le C_*,
 \quad C_*:=\frac{J_s^2}{\lambda_-}e^{2\lambda_+}.
\end{equation}
Impose finitely many site requirements $T_I\subset I,T_J\subset J$.
On occupation $O$, $|O|=m$, the diagonal of the column covariance
vanishes unless $T_J\subset O$. Otherwise, with
$K=T_J\cup(T_I\cap O)$, it equals
\begin{equation}
 w_N\binom{m-|K|}{p-|K|}
       \binom{N-m+p-|T_I|}{p-|T_I|}.
\end{equation}
Invalid binomial coefficients are zero. Thus, on an interior charge
interval, with $t=|T_I|+|T_J|$ fixed, the row and column covariances
are bounded by
\begin{equation}\label{restricted}
 C_m\frac{(p)_{|T_J|}}{(m)_{|T_J|}}
       \frac{(p)_{|T_I|}}{(N-m+p)_{|T_I|}}
 \le C_* (K_t p/N)^t.
\end{equation}
The row version swaps the two label roles.

For a fixed Gaussian pairing, scan the matrix word from right to left.
At the first endpoint of an index use the column map
$V\psi=\sum_\alpha A_\alpha\psi\otimes e_\alpha$; at its second
endpoint use the row map. Their norms are the square roots of the
corresponding covariance norms. Open indices occupy distinct ancillas;
crossings only permute these ancillas unitarily. Fixed CAR sandwiches
are contractions. The complete pairing contraction is therefore
bounded by the product of endpoint norms. Requirements at both
endpoints are first combined into the same allowed index family.
This argument retains the CAR signs until the matrix contraction is
bounded. Taking absolute local weights first would lose a factor
exponential in the square of the degree.

\newpage
\section{Pure marks, repeated deficits and Gaussian projection}
The exact body product rule is
\begin{equation}
 D(XY)=(DX)Y+X(DY)-\sum_i[c_i^\dagger,X][Y,c_i].
\end{equation}
Since $DA_{IJ}=pA_{IJ}$, a deficit $(np-D)$ on a product of $n$
slots selects two slots and one site. The two commutators require
$i\in J_a$ and $i\in I_b$. In a pure Wick block those slots cannot
pair with each other, so these are two independent set labels.
The single-slot part of $p-\LL$ requires $i\in I\cap J$ and has
the same property. For $s$ distinct sites,
\begin{equation}
 N^s(p/N)^{2s}=\lambda_N^s.
\end{equation}
Combining this count with \eqref{restricted}, a word with $L$ slots
and $k$ fixed pure marks satisfies
\begin{equation}\label{word}
 |\tau_N(\text{marked word})|
 \le K_k(1+L)^{2k}C_*^{L/2}(L-1)!!.
\end{equation}
The outside-charge estimate uses at worst $N^k$ and the binomial
tail; its exponential decay absorbs this fixed power. For four
copies of a one-mark degree-$n$ field,
$(4n-1)!!\le16^n(n!)^2$, giving \eqref{envelope}.

Repeated deficits on one pure block also obey this bound. Extend
the body operator to a homogeneous parity-$a$ word by
\begin{equation}
 D_gX=\tfrac12\{Q,X\}-(-1)^a\sum_i c_i^\dagger Xc_i.
\end{equation}
A normal ordered word with $r_c+r_a$ letters has eigenvalue
$(r_c+r_a)/2$. Each contraction removes two letters. Applying the
next deficit either counts the already removed pairs or adds another
tagged contraction. The multiplicity depends on the fixed number of
marks, with no new $p$ loss. Therefore any fixed sequence of $k$
$D/\LL$ deficits has bound
$K_{k,r}(1+n)^{2k}L_{k,r}^n\sqrt{n!}$.

There is also a dimension independent, same-norm Gaussian tool.
For $\deg P\le n$, the contraction
$\mathsf O_tP=\sum_{j\le n}t^j\WW_jP$, $-1\le t\le1$, has a
Chebyshev expansion whose coefficients have norm at most $2\|P\|_r$.
Leading coefficients and absolute coefficient sums give
\begin{align}\label{topchaos}
 \|\WW_nP\|_r&\le2^n\|P\|_r,\\
 \left\|\sum_{j\le n}\theta_j\WW_jP\right\|_r
 &\le K(1+\sqrt2)^n\max_j|\theta_j|\,\|P\|_r.\label{chaos}
\end{align}
Linear Gaussian matrix moments \cite{Tropp} yield
$\|H_m^n\|_r\le K_r^n\sqrt{n!}$; hence the same envelope holds
for $Z_n$. No matrix-valued $S_2$ to $S_4$ hypercontractivity is assumed.

\newpage
\section{The covariance is an exact completely positive path}
On $\mathcal H_k=\wedge^k\mathbb C^N$ define
\begin{equation}
 T_kX=\frac1k\sum_i c_i^\dagger Xc_i,
 \qquad S_kX=\frac1{N-k}\sum_i c_iXc_i^\dagger.
\end{equation}
$T_k$ maps charge $k-1$ to $k$ and $S_k$ maps $k+1$ to $k$.
Both are unital and preserve their respective normalized sector
traces. Thus they contract every Schatten norm, also after disorder
averaging. The deterministic covariance map has the exact factorization
\begin{equation}\label{path}
 \Phi_m(X)=w_N\sum_{I,J} C_I^\dagger C_JX C_J^\dagger C_I,
 \qquad
 \frac{\Phi_m}{C_m}
 =T_m\cdots T_{m-p+1}S_{m-p}\cdots S_{m-1}.
\end{equation}
Expanding the Kraus products proves this identity: each unordered
$p$-set has $p!$ orders on each side, the reordering signs cancel,
and the path denominators give
$(m)_p(N-m+p)_p=(p!)^2\binom mp\binom{N-m+p}p$.

Put $\bar D=2\LL-D$. For any neutral full-Fock section $X$,
\begin{align}\label{step}
 T_kX_{k-1}-(1-b/k)X_k&=((b-D)X)_k/k,\\
 S_kX_{k+1}-(1-b/(N-k))X_k&=((b-\bar D)X)_k/(N-k).
\end{align}
Set
\begin{align}
 \theta_{v,m}&=\prod_{a=0}^{p-1}
  (1-vp/(m-a))_+(1-vp/(N-m+p-a))_+,\\
 A_{v,m}&=p(\Phi_m-C_m\theta_{v,m})Z_{v,m}.
\end{align}
Inactive sectors are defined separately by zero covariance. In
$N/3\le m\le2N/3$, the path stays in a wider interior interval.
If $vp\le N/8$, telescope \eqref{step} over $2p$ steps. Every
source has coefficient at most $Kp/N$ after multiplying by $pC_m$;
their sum costs $p^2/N=\lambda_N$. The pure marked bound therefore
proves \eqref{envelope} for $A_v$. If $vp>N/8$, then
$p<8\lambda_Nv$, and the coarse CP estimate
$\|A_v\|_r\le2pC_*\|Z_v\|_r$ suffices. Outside the central
interval the same estimate is absorbed by the binomial tail.

This supplies a bound for the actual covariance residual, before
taking any degree or size limit. A coefficientwise occupation
Markov estimate alone would not control its summed Gaussian
Schatten norm; the full CP factorization is the needed bridge.

\newpage
\section{The actual neutral time series}
Retain finite-sector parameters in the Jacobi recurrence:
\begin{equation}
 a_{j,m}=\sum_{v<j}\theta_{v,m},\qquad
 Q_{j+1}(x)=xQ_j(x)-C_ma_{j,m}Q_{j-1}(x),
 \qquad \eta_j=p[Z_j-Q_j(H)].
\end{equation}
The exact Gaussian creation/annihilation identity from the previous
neutral stage \cite{Neutral} is
\begin{equation}\label{eta}
 pr_j=\sum_{v+w=j-1}A_v\diamond Z_w,
 \qquad \eta_{j+1}=H\eta_j-C_ma_j\eta_{j-1}-pr_j.
\end{equation}
Since $A_v$ is pure coupling degree $v$,
\begin{equation}
 A_v\diamond Z_w=\WW_{v+w}(A_vH^w).
\end{equation}
Equation~\eqref{topchaos} and fixed $2r$ H\"older give an exponential
times $\sqrt{(v+w)!}$ envelope. Unroll \eqref{eta} before taking
norms. Its terms are bounded Jacobi path coefficients times
$H^a\WW_{v+w}(A_vH^w)$. Only a fixed $4r$ input is needed,
independent of the degree. The path count is exponential and
$\sqrt{a!v!w!}\le\sqrt{j!}$, giving \eqref{envelope} for $\eta_j$.
Outer charges can instead use $p(Z_j-Q_j(H))$ and the binomial tail.

The exact finite-parameter expansion
$e^{isH_m}=\sum_jc_{j,m}(s)Q_{j,m}(H_m)$ obeys
$|c_{j,m}(s)|\le K_S L_S^j/j!$. Subtracting a function of $H_m$
and using exact pinching gives
\begin{align}\label{neutral}
 p(1-\PP)\Phi_m(e^{isH_m})
 &=(1-\PP)\sum_{j\ge0}c_{j,m}(s)T_{j,m},\\
 T_{j,m}&=A_{j,m}+(C_m\theta_{j,m}-\Phi_m)\eta_{j,m}.
\end{align}
The normalized covariance is a contraction, so $T_j$ satisfies
\eqref{envelope}. With $\alpha_m=1-p/m$ and
$\epsilon_m=1-\alpha_m^2$, the actual neutral memory is
\begin{equation}
 B_{m,N}(s)=\frac{m\epsilon_m}{p}(1-\PP)e^{-iH_ms/2}
 \sum_jc_{j,m}(\alpha_ms)T_{j,m}e^{-iH_ms/2}.
\end{equation}
Here $0\le m\epsilon_m/p\le2$ on active sectors, and the phases
are unitary. Thus the actual remainder, uniformly at finite $N$, is
bounded by
\begin{equation}\label{tail}
 K_{r,S}\sum_{j>K}(j+1)^{b_r}L_{r,S}^j/\sqrt{j!}
 \longrightarrow0.
\end{equation}
The previously constructed local-size/ordered-Wick proxy has the
same bound by Section~3. The old compact-time $L^2$ strong bridge
therefore upgrades by interpolation to every fixed $S_r$.

\newpage
\section{Differentiating the CP path without a size loss}
Use distinct notation for unnormalized maps
$K_-X=\sum c_i^\dagger Xc_i$, $K_+X=\sum c_iXc_i^\dagger$,
and normalized maps $t=Q^{-1}K_-$, $s=(N-Q)^{-1}K_+$.
Direct CAR give
\begin{align}\label{comm}
 [D,t]&=t-t^2,&[\bar D,s]&=s-s^2,\\
 [D,s]&=-s+\frac Q{N-Q}ts-\frac{2Q-N}{N-Q}I,\\
 [\bar D,t]&=-t+\frac{N-Q}{Q}st+\frac{2Q-N}{Q}I.
\end{align}
For a scalar charge function, retain the finite differences:
\begin{align}\label{diff}
 [D,f(Q)]&=[f(Q)-f(Q-1)]K_-,\\
 [\bar D,f(Q)]&=[f(Q)-f(Q+1)]K_+.
\end{align}
Commuting $D/p$ or $\LL/p$ through at most $2p$ normalized
steps costs only a bounded factor on the interior interval.
Apply this to the telescope for $A_v$, not to its coarse difference.
Each term is $pC(Q)f_h(Q)P_hE_h$, with a pure deficit $E_h$,
$|f_h|\le K/N$, and
$|\Delta f_h|\le K(1+v\lambda_N)/N^2$.
Repeated pure marks control derivatives of $E_h$; \eqref{diff}
controls all scalar coefficients. For $vp>N/8$, use
$p\le8\lambda_+v$ and $N\le64\lambda_+v^2$ instead.
Outer charges are controlled before applying any full-trace theorem.
The result is
\begin{equation}\label{Ader}
 \|A_v\|_r+\|DA_v/p\|_r+\|\LL A_v/p\|_r
 \le K_r(v+1)^{b_r}L_r^v\sqrt{v!}.
\end{equation}

The Clifford Riesz theorem of Ben Efraim--Lust-Piquard,
Theorem 6.8(1),(3), and its semigroup interpolation Lemma 5.6
\cite{Riesz}, supply the product bridge. In the present convention
$N_{\rm Clifford}=2\LL$. For either physical product-rule cross
term $\Gamma_\sigma(X,Y)$, $\sigma=D,\LL$,
\begin{equation}\label{riesz}
 \|\Gamma_\sigma(X,Y)/p\|_r
 \le K_r\bigl[\|X\|_{2r}\|\LL X/p\|_{2r}
                   \|Y\|_{2r}\|\LL Y/p\|_{2r}\bigr]^{1/2}.
\end{equation}
These are normalized annealed \emph{full-Fock} norms. The theorem
is not applied within a fixed charge sector. Its row and column
square functions must both be kept for non-Hermitian inputs.

\newpage
\section{Ordinary derivatives and the Jacobi error}
A useful independent CP estimate avoids auxiliary dimension factors.
If $\Phi$ is completely positive, with trace adjoint $\Phi^\dagger$,
then scalar H\"older on the spectral decompositions of positive $X,Y$
gives, for conjugate $u,v$,
\begin{equation}
 \tr Y\Phi(X)\le
 [\tr X^u\Phi^\dagger(I)]^{1/u}
 [\tr Y^v\Phi(I)]^{1/v}.
\end{equation}
Disorder can be included in that same H\"older inequality. Duality
with $u=2r$ then gives
\begin{equation}
 \|\Phi(X)\|_r\le
 \|X\|_{4r}\|\Phi^\dagger(I)\|_2^{1/(2r)}
                 \|\Phi(I)\|_{2r-1}^{1-1/(2r)}.
\end{equation}
For arbitrary $X$ an absolute factor at most four suffices.
For the force map
\begin{equation}
 \Gamma_H(X)=\sum_i[c_i^\dagger,H]X[H,c_i],
\end{equation}
the Clifford square-function bound controls both Grams by $O_r(p)$.
Consequently $\|\Gamma_H(X)\|_r\le K_rp\|X\|_{4r}$.
The exact ordered identity
\begin{equation}
 DH^w=pwH^w-
 \sum_{a+b+c=w-2}H^a\Gamma_H(H^b)H^c
\end{equation}
and the analogous $\LL$ formula imply
\begin{equation}\label{ordinaryder}
 \|H^w\|_r+\|DH^w/p\|_r+\|\LL H^w/p\|_r
 \le K_r(w+1)^{b_r}L_r^w\sqrt{w!}.
\end{equation}
Alternatively, the tagged proof gives this independently: ordinary
internal pairs may leave only one independent set requirement per
site, but each normalized deficit supplies $1/p$. For $s\le k$
distinct sites in $k$ marks, $p^{-k}N^s(p/N)^s=p^{s-k}\le1$.

Now unroll \eqref{eta}. Its terms have the form
\begin{equation}
 b(Q)H^a\WW_{v+w}(A_vH^w).
\end{equation}
Use \eqref{riesz} once inside the projection and once outside it;
physical derivatives commute with Gaussian projection. The needed
norm order is a fixed multiple of $r$, not one that grows at each
recursive step. Coefficient derivatives use \eqref{diff} separately.
The central Jacobi path weights and their normalized differences
grow at most exponentially times a degree polynomial. The outer
estimate follows directly from $p(Z_j-Q_j(H))$ and its derivatives.
It follows that
\begin{equation}\label{etader}
 \|\eta_j\|_r+\|D\eta_j/p\|_r+\|\LL\eta_j/p\|_r
 \le K_r(j+1)^{b_r}L_r^j\sqrt{j!}.
\end{equation}

\newpage
\section{The ordinary charged tail: keep the charge differences}
The inverse Jacobi expansion is exact:
\begin{equation}
 H^n=\sum_{j=0}^nh_{nj}(Q)[Z_j-\eta_j/p].
\end{equation}
Its coefficients are weighted up/down paths. On an interior charge
interval the recurrence weights $\beta_j=C_m\sum_{v<j}\theta_v$
are bounded. The sharper difference estimate needed here follows
from
\begin{equation}
 \frac{C_m}{C_{m-1}}=
 \frac m{m-p}\frac{N-m+1}{N-m+p+1},\qquad
 N|\Delta C_m|\le K\lambda_N(1+|X_m|),
 \quad X_m=\frac{m-N/2}{p}.
\end{equation}
Each positive-part factor in $\theta_v$ is Lipschitz, including
across its cutoff. All central denominators are at least $cN$, so
\begin{equation}
 N|\Delta\theta_{v,m}|\le K\lambda_Nv
 \quad\text{for every }v.
\end{equation}
Taking one difference along each path, rather than differentiating
every coefficient simultaneously, gives
\begin{equation}\label{hdiff}
 \sum_j|h_{nj,m}|\le K^n,\qquad
 \sum_jN|\Delta h_{nj,m}|
 \le K(n+1)^3L^n(1+|X_m|).
\end{equation}
Binomial full-Fock charge has bounded fixed moments of $X_m$.

Define the actual ordinary charged coefficient
\begin{equation}
 G_{n,N}=[p\NN-D-\lambda_N\NN(\NN-1)]H^n.
\end{equation}
With $\Delta h_m=h_m-h_{m-1}$, its non-counterterm part is
\begin{align}\label{reassembly}
 (p\NN-D)H^n=\sum_j\bigl[&h_{nj}(pj-D)Z_j
       -(\Delta h_{nj})K_-Z_j-h_{nj}\NN\eta_j\notag\\
       &+h_{nj}D\eta_j/p
       +(\Delta h_{nj})K_-\eta_j/p\bigr].
\end{align}
This retains all internally contracted Gaussian terms. The first
term is pure marked; the second uses \eqref{hdiff}; the third uses
\eqref{chaos}; the last two use \eqref{etader} and coefficient
differences. On outer charges use the original $G_n$ and the coarse
$D=O(N)$ bound with exponential binomial tails. Finally
$\lambda_N\NN(\NN-1)$ is controlled by \eqref{chaos}.
Thus \eqref{envelope} holds for $G_n$ and for
\begin{equation}
 K_n=G_{n+2}-\{H,G_{n+1}\}+HG_nH.
\end{equation}
The actual series $\sum_n(is)^nG_n/n!$ and its fixed derivatives
and energy insertions therefore have the finite-$N$ tail
\eqref{tail}. This supplies the missing ordinary, rather than
only pure-Wick, bound.

\newpage
\section{Returning to the actual fields}
The preceding memory reduction \cite{Transfer,Joint} is also valid
in every fixed $S_r$ norm. Here is the required upgrade, specifying
which bounds are used.

First, the charged OU remainder is exactly
\begin{equation}
 E_m-G_N=R_m(\NN)e^{isH},\qquad
 E_m=[m(1-\alpha_m^{\NN})-D]e^{isH}|_m.
\end{equation}
The physical derivative cancels. The central bound
$|R_m(k)|\le a_m(k+1)^3$, with
$\|a_m\|_{L^r(\mathrm{binomial})}=O_r(p^{-1})$, combines with
\eqref{chaos} and the ordinary moment envelope to give $O(p^{-1})$
for fixed derivatives and energy insertions. The other $\alpha_m-1$
terms vanish by the new charged bounds (or by the earlier coarse
high-moment interpolation).

Second, the old positive-force estimates bound $p\Lambda(s)$ in
every fixed norm. Its $L^2$ convergence to the energy regression
therefore upgrades by interpolation. Fixed-order H\"older transfers
this to the first return; the next return is $O(p^{-1})$.
The four-auxiliary-Gaussian energy regression and the next OU
remainder are treated at fixed higher norm orders in the same way.
This gives the actual memory to charged-plus-neutral strong bridge.

Third, for each fixed $n$, the joint-regression result uses the same
original disorder:
\begin{equation}\label{Rbridge}
 R_{n,N}:=G_{n,N}-g_n(X_N,H_N)
 \longrightarrow \sum_d a_{nd}J_{d,N}\quad\text{strongly in }L^2.
\end{equation}
The new $G_n$ bounds, the fixed-degree charge/energy polynomial
bounds, and pure marked bounds upgrade this to $S_r$.
Ordinary internal pairs are removed through \eqref{Rbridge}; they
are not silently treated as forbidden pure-Wick pairs.

The neutral local proxy is a finite sum of ordered pure marked
blocks and ordinary outside powers at each fixed degree. Its
previous strong bridge and the new two-sided high bounds give the
same upgrade. Fixed residuals have no diagonal spectral mass by the
preceding two-point argument. Exact pinching first vanishes in
$L^2$, then in $S_r$ by interpolation and contraction. Only then
are the finite-degree common-source moments used.

Here $\FF_N(s)$ is the actual centered memory field:
$m(1-\PP)e^{-iH_ms/2}M_m(s)e^{-iH_ms/2}$, with the compressed-bath
memory $M_m$ of the preceding reconstruction. The established
bilocal return estimate \cite{Four} is
\begin{align}\label{return}
 \sup_{|t|\le S}\|\BB_N(t)-\widetilde\BB_N(t)\|_r
 &\le K_{r,S}p^{-1/(2r-1)}+K_{r,S}pe^{-N/(8r)},\\
 \widetilde\BB_N(t)&=-\int_0^t ds
 \int_{-(t-s)/2}^{(t-s)/2}du\,
 e^{iH_Nu}\FF_N(s)e^{-iH_Nu},\quad t\ge0.\label{filter}
\end{align}
At $r=4$ the return error is $O(p^{-1/7})$.

\newpage
\section{A constructive complete four-point function}\label{construction}
Let $\tau_\infty$ be the limiting common-source tracial functional.
At fixed degree it is explicitly evaluated as follows. Expand all
fields and counterterms into ordinary outside powers and pure marked
blocks; keep the whole-block Wick restrictions and the true matrix
order. For every allowed pairing, form the local two-label CAR
polynomial $T_{uv}(z)$ and divide by its source-dependent singleton
factors, obtaining $R_{uv}(z)$. Extract the required nilpotent source
coefficient from
\begin{equation}\label{pairing}
 (Ce^{-\lambda})^\ell
 \exp\left\{\lambda\sum_{u<v}[R_{uv}(z)-1]\right\},
\end{equation}
then sum the allowed pairings. Counterterms and chaos projectors
are expanded with their actual coefficients. The prefactor follows
from $w_N\binom Np^2 2^{-2p}\to Ce^{-\lambda}$ per chord.
Finite-size overlap exponential moments justify this fixed-word
Poisson limit \cite{Four}; Section~3, separately, gives summable
degree bounds. General descendants are not replaced by the
three-species low-field mixed-$q$ law.

Write $R_n^\infty=\sum_d a_{nd}J_d$ and
$L_n=R_{n+2}^\infty-\{H,R_{n+1}^\infty\}+HR_n^\infty H$.
Let $N_j$ be the centered neutral ordered descendant. For
$q=e^{-4\lambda}$, let $H_j$ be the monic $q$-Hermite polynomial
of squared norm $C^j[j]_q!$ in $\rho_q$, and put
$c_j(s)=\int e^{isx}H_j(x)d\rho_q(x)/(C^j[j]_q!)$. Define
\begin{equation}
 \FF^{[K]}(s)=e^{-iHs/2}\left[
 \sum_{n=0}^K\frac{(is)^n}{n!}L_n
       +2C\sum_{j=1}^Kc_j(s)N_j\right]e^{-iHs/2}.
\end{equation}
The neutral sign is positive. Apply the filter in \eqref{filter}
to define $\BB^{[K]}(t)$, and extend to negative times by adjoints.
Fixed-degree strong replacement and $\|1-\PP\|_r\le2$ make the
limiting centered blocks inherit the degree envelope; this does not
rely on $L^2$ regression being an $S_r$ contraction.
For $\varepsilon_i\in\{1,*\}$, center times $T_i$, and relative
times $t_i$, the answer is
\begin{equation}\label{four}
 \boxed{\lim_{K\to\infty}\tau_\infty
 \prod_{i=1}^4
 \left(e^{iHT_i}\BB^{[K]}(t_i)e^{-iHT_i}\right)^{\varepsilon_i}.}
\end{equation}
At fixed $K$, expand the energy exponentials into polynomials and
evaluate them by \eqref{pairing}; Gaussian energy moments control
this auxiliary truncation. Thus \eqref{four} is a convergent
algorithm, without requiring a closed four-field spectral kernel.

For completeness, the filter norm is at most $S^2/2$. The new
$S_4$ tails make its truncated fields Cauchy, and all actual fields
are bounded in $S_4$. Telescoping a four-factor product and using
H\"older bounds its error by one $S_4$ difference times three
$S_4$ bounds. First pass $N\to\infty$ at fixed truncation using
the strong interfaces; then remove truncation using \eqref{tail}.
The return contribution is controlled by \eqref{return}. All these
steps are locally uniform in the fixed time windows. Unitary
insertions preserve strong-error bounds for arbitrary center times,
but spectral convergence itself is only asserted on compact windows.

\newpage
\section{Checks, provenance and the next physical step}
Three independent finite algebra checks accompany the proof. They
test exact identities and nontrivial CAR signs; they do not substitute
for the all-degree estimates.
\begin{itemize}
\item Restricted covariance and ancilla contraction:
42,336 covariance comparisons, 564 networks (452 crossed),
94,743 label assignments, and 1,776 exact rational sector bounds.
There are 113 nonzero crossed CAR-sandwich cases.
\item Coherence/CP channel:
340 coherence identities, 265 normalized path identities,
75 inactive cases, 2,609 kernel entries and 1,305 tagged-site bounds,
including 132 nonconstant inputs and 14 nontrivial CAR gauges.
\item Charged reassembly: 3,312 matrix identities, 9,949 entry
comparisons and 84 polynomial identities, including 16 nonzero
charge-difference fixtures. The polynomial fixtures use one Gaussian
coordinate to check algebra, not the full ensemble's moments.
\end{itemize}
Drivers and outputs are stored under \texttt{code/} and
\texttt{results/susskind-bilocal-fourth-tail/}. The local evidence
manifest binds the input/output hashes, parameters, dependency
versions, parent manifests, build logs and final page renders.
The two preceding checks were first successful runs; old scientific
outputs and failed records are preserved. The new third check is
also saved independently, rather than overwriting a predecessor.

The theoretical increment is precise: complete finite-relative-time
four-point functions now follow from actual finite-size uniform
tails, rather than from a formal limiting Taylor series. The next
question is which state and polarization can turn this tracial
collective algebra into a physical particle construction, and how
the microscopic Hamiltonian supplies spatial/internal generators.
The positive low-channel commutator remains a useful calibration.
It is not by itself a finite-$G$ gravitational scattering amplitude.
No new issue reply, Git commit or push accompanies this local stage.

\begin{thebibliography}{9}\setlength{\itemsep}{2pt}
\bibitem{Riesz} L. Ben Efraim and F. Lust-Piquard,
\emph{Poincar\'e type inequalities on the discrete cube and in the
CAR algebra}, Probab. Theory Relat. Fields \textbf{141} (2008),
569--602, \href{https://arxiv.org/abs/math/0702233}{math/0702233}.
Theorem 6.8(1),(3) and Lemma 5.6.
\bibitem{Tropp} J. A. Tropp,
\emph{User-Friendly Tail Bounds for Sums of Random Matrices},
Found. Comput. Math. \textbf{12} (2012), 389--434,
\href{https://arxiv.org/abs/1004.4389}{arXiv:1004.4389}, Sec. 4.7.
\bibitem{Four} A. Day and Chengheng, \emph{A common-disorder algebra
for two bilocal channels}, local draft (7 Nov. 2026), with its
multisource and actual-transfer proof notes.
\bibitem{Neutral} A. Day and Chengheng, \emph{The neutral marked
series}, local draft and proof notes (7 Nov. 2026).
\bibitem{Transfer} A. Day and Chengheng, \emph{The oriented memory
transfer} and \emph{Two-component memory reduction}, local drafts
and proof notes (7 Nov. 2026).
\bibitem{Joint} A. Day and Chengheng, \emph{Joint charge/energy
regression and the complete memory kernel}, local draft and proof
notes (7 Nov. 2026).
\bibitem{Recon} A. Day and Chengheng, \emph{Strong reconstruction
of the complete bilocal fluctuation}, local draft (7 Nov. 2026).
\end{thebibliography}
\end{document}
